% UAS Logika Predikat --- soal uraian (format mengikuti UTS/uts.tex).
% Soal diadaptasi dari sumber daring; skenario berbeda dari tugas/tugas2.tex.
\documentclass[12pt,a4paper]{article}

\usepackage[utf8]{inputenc}
\usepackage[T1]{fontenc}
\usepackage[indonesian]{babel}

\usepackage{geometry}
\geometry{margin=2cm}

\usepackage{lmodern}
\usepackage{amsmath,amssymb}
\usepackage{enumitem}
\usepackage{tabularx}
\usepackage{booktabs}
\usepackage{microtype}

\usepackage{fancyhdr}
\usepackage{lastpage}

\setlength{\headheight}{14pt}
\addtolength{\topmargin}{-2pt}
\usepackage[hidelinks,breaklinks]{hyperref}

\newcommand{\JudulUjian}{Ujian Akhir Semester (UAS)}
\newcommand{\KodeMK}{TIF-xxx}
\newcommand{\NamaMK}{Logika Matematika}
\newcommand{\Semester}{Gasal 2025/2026}
\newcommand{\MateriUjian}{Logika Predikat (kalkulus predikat)}
\newcommand{\WaktuUjian}{120 menit}
\newcommand{\NilaiMaksimal}{100}
\newcommand{\Prodi}{Teknik Informatika}
\newcommand{\Fakultas}{Fakultas Teknologi Informasi}
\newcommand{\Institusi}{Universitas Bale Bandung}
\newcommand{\DosenPengampu}{[Nama Dosen Pengampu]}
\newcommand{\TanggalUjian}{[DD Bulan YYYY]}

\pagestyle{fancy}
\fancyhf{}
\fancyhead[L]{\small UAS --- \NamaMK}
\fancyhead[R]{\small \Semester}
\fancyfoot[C]{\small Halaman \thepage\ dari \pageref*{LastPage}}
\renewcommand{\headrulewidth}{0.4pt}
\renewcommand{\footrulewidth}{0pt}

\setlength{\parindent}{0pt}
\setlength{\parskip}{0.5em}

\newcounter{soaluas}
\newcommand{\soaluas}{%
  \refstepcounter{soaluas}%
  \par\vspace{0.8ex plus 0.15ex}%
  \noindent\textbf{Soal \thesoaluas.}\quad
}

\newcommand{\ruangjawaban}[1][6cm]{%
  \par\vspace{0.4em}%
  \noindent\rule{\linewidth}{0.4pt}\\[0.35em]%
  \rule{\linewidth}{0.4pt}\\[0.35em]%
  \ifdim#1>6cm\rule{\linewidth}{0.4pt}\\[0.35em]\fi
  \ifdim#1>8cm\rule{\linewidth}{0.4pt}\\[0.35em]\fi
  \par\vspace{0.7em}%
}

\begin{document}

\begin{center}
  {\Large\bfseries \JudulUjian}\\[0.35em]
  {\large \NamaMK\ (\KodeMK)}\\[0.25em]
  {\normalsize \Semester}\\[0.8em]
\end{center}

\noindent
\begin{tabularx}{\linewidth}{@{}p{4.8cm} @{\hspace{0.35em}:\hspace{0.65em}} >{\raggedright\arraybackslash}X@{}}
  Nama Mahasiswa        & \rule{0.72\linewidth}{0.4pt} \\
  Nomor Induk Mahasiswa & \rule{0.72\linewidth}{0.4pt} \\
  Kelas / Kelompok      & \rule{0.72\linewidth}{0.4pt} \\
  Program Studi         & \Prodi \\
  Fakultas              & \Fakultas \\
  Institusi             & \Institusi \\
  Dosen Pengampu        & \DosenPengampu \\
  Tanggal ujian         & \TanggalUjian \\
  Waktu pengerjaan      & \WaktuUjian \\
  Nilai maksimal        & \NilaiMaksimal \\
  Cakupan materi        & \MateriUjian \\
\end{tabularx}

\vspace{0.8em}
\hrule
\vspace{0.8em}

\section*{Petunjuk Umum}
\begin{itemize}[leftmargin=*, itemsep=0.25em]
  \item Ujian terdiri atas \textbf{10 soal uraian} (Soal 1--10) pada materi \textit{Logika Predikat}.
        Setiap soal memuat label bahasan (cetak miring); jawab sesuai butir \textbf{a, b, c, \ldots} bila ada.
  \item Isi \textbf{nama, NIM, dan kelas} pada bagian identitas sebelum mulai.
        Tulis jawaban \textbf{berurutan per nomor soal} pada garis yang disediakan; bila tidak cukup, lanjutkan di lembar tambahan yang dilekatkan dan beri penanda ``Soal X''.
  \item Gunakan predikat, konstanta, fungsi, variabel, serta kuantor \(\forall\) dan \(\exists\) secara konsisten.
        Penghubung logik: \(\neg\), \(\wedge\), \(\vee\), \(\rightarrow\), \(\leftrightarrow\).
        Nilai kebenaran boleh ditulis \(T/F\) atau Benar/Salah.
  \item Definisikan predikat/fungsi/konstanta yang Anda pakai sebelum menuliskan formula, kecuali soal sudah menyediakan kunci terjemahan.
  \item Ujian bersifat \textbf{individual}; dilarang bekerja sama, menyalin, atau menggunakan perangkat elektronik serta catatan kecuali diizinkan tertulis pengawas.
  \item Alokasi waktu disarankan sekitar 10--12 menit per soal dalam \WaktuUjian; prioritaskan menyelesaikan semua nomor.
  \item Lembar ujian \textbf{wajib dikumpulkan} kepada pengawas segera setelah waktu \WaktuUjian\ habis.
\end{itemize}

\section*{Soal Uraian}

\soaluas
\textit{(Pendahuluan / Motivasi)}\\
\begin{enumerate}[label=\alph*., leftmargin=2em, itemsep=0.25em]
  \item Jelaskan mengapa kalkulus proposisi \textbf{tidak memadai} untuk menyatakan kalimat seperti
        ``Semua mahasiswa TI wajib mengambil mata kuliah logika'' atau
        ``Ada dosen yang mengajar dua kelas sekaligus''.
        Sebutkan peran \textbf{objek} dan \textbf{kuantor} dalam jawaban Anda.
  \item Berikan \textbf{satu} contoh predikat unary dan \textbf{satu} contoh predikat biner dari konteks kampus,
        lalu tuliskan aritas masing-masing.
\end{enumerate}
\ruangjawaban[8cm]

\soaluas
\textit{(Terjemahan ke Formula)}\\
Terjemahkan pernyataan berikut ke logika predikat.
Definisikan predikat (dan konstanta bila perlu) yang Anda gunakan.
\begin{enumerate}[label=\alph*., leftmargin=2em, itemsep=0.2em]
  \item Ada orang yang merupakan kakek/nenek dari dirinya sendiri
        (gunakan \(\mathit{Person}(p)\) dan \(\mathit{ParentOf}(p_1,p_2)\):
        \(p_1\) adalah orang tua dari \(p_2\)).
  \item Hanya asisten laboratorium yang boleh membuka ruang server.
  \item Tidak ada orang yang mengagumi siapa pun
        (domain: orang; gunakan \(\mathit{Admires}(x,y)\)).
  \item Tidak ada yang berbentuk limas kecuali benda itu kecil
        (pola ``nothing \ldots unless \ldots''; ganti konteks: \(\mathit{Pyramid}(x)\), \(\mathit{Small}(x)\)).
  \item Politisi dapat menipu sebagian rakyat sepanjang waktu,
        dan dapat menipu seluruh rakyat pada sebagian waktu,
        tetapi tidak dapat menipu seluruh rakyat sepanjang waktu
        (gunakan \(\mathit{Politician}(x)\), \(\mathit{Citizen}(y)\), \(\mathit{Time}(t)\), \(\mathit{Fools}(x,y,t)\)).
\end{enumerate}
\ruangjawaban[9cm]

\soaluas
\textit{(Formula ke Bahasa Indonesia)}\\
Ubahlah formula berikut menjadi pernyataan bahasa Indonesia yang alami (hindari menyebut nama variabel secara eksplisit).
Gunakan kunci:
\(M(x)\): \(x\) adalah mahasiswa,
\(D(x)\): \(x\) adalah dosen,
\(C(y)\): \(y\) adalah mata kuliah,
\(T(x,y)\): \(x\) mengambil mata kuliah \(y\),
\(A(x,y)\): \(x\) mengajukan pertanyaan kepada \(y\),
\(L(x,y)\): \(x\) lebih tinggi nilainya daripada \(y\).
\begin{enumerate}[label=\alph*., leftmargin=2em, itemsep=0.25em]
  \item \(\forall x\,(M(x) \rightarrow \exists y\,(C(y) \wedge T(x,y)))\)
  \item \(\exists x\,\bigl(M(x) \wedge \forall y\,((M(y) \wedge \neg(x=y)) \rightarrow L(x,y))\bigr)\)
  \item \(\forall x\,\bigl(M(x) \rightarrow \exists y\,(D(y) \wedge A(x,y))\bigr)\)
  \item \(\forall x\,\forall y\,\bigl((\mathit{Adjoins}(x,y) \wedge \mathit{Cube}(y)) \rightarrow \mathit{Tet}(x)\bigr)\)
        (boleh diterjemahkan bebas: ``hanya tetrahedron yang bersebelahan dengan kubus'' atau padanan yang setara).
\end{enumerate}
\ruangjawaban[8cm]

\soaluas
\textit{(Negasi Kuantor)}\\
Untuk setiap pernyataan berikut:
(i)~tuliskan formula logika predikat,
(ii)~negasikan formula tersebut dan dorong \(\neg\) sedalam mungkin,
(iii)~terjemahkan hasil negasi kembali ke bahasa Indonesia.
\begin{enumerate}[label=\alph*., leftmargin=2em, itemsep=0.25em]
  \item Semua dosen di program studi ini bergelar doktor.
  \item Ada mahasiswa yang lulus semua mata kuliah wajib.
  \item Tidak ada asisten yang terlambat lebih dari dua kali.
\end{enumerate}
\ruangjawaban[9cm]

\soaluas
\textit{(Term, Proposisi, Kalimat)}\\
Asumsikan: \(c,d\) konstanta; \(u,v\) variabel; \(h\) fungsi unary; \(k\) fungsi biner; \(R\) predikat biner; \(S\) predikat unary.
Untuk setiap ekspresi, tentukan apakah ia \textbf{term}, \textbf{proposisi}, \textbf{kalimat}, atau \textbf{bukan ekspresi yang terbentuk dengan baik}.
Berikan alasan singkat.
\begin{enumerate}[label=\alph*., leftmargin=2em, itemsep=0.2em]
  \item \(k(c,h(d))\)
  \item \(R(h(u),c)\)
  \item \(\forall u\,S(u)\)
  \item \(S(R(c,d))\)
  \item \(\mathit{if}\; S(c)\; \mathit{then}\; h(d)\; \mathit{else}\; k(c,d)\)
  \item \(R(c) \vee S(d)\)
\end{enumerate}
\ruangjawaban[8cm]

\soaluas
\textit{(Variabel Bebas / Terikat \& Kalimat Tertutup)}\\
Untuk setiap formula, tandai setiap kemunculan variabel sebagai \textbf{bebas} atau \textbf{terikat}.
Jika terikat, sebutkan kuantor yang mengikatnya.
Nyatakan pula apakah formula merupakan \textbf{kalimat tertutup}.
Perhatikan cakupan tanda kurung.
\begin{enumerate}[label=\alph*., leftmargin=2em, itemsep=0.25em]
  \item \((\forall u)\,R(u,v) \wedge S(u)\)
  \item \((\forall u)\,(R(u,v) \wedge S(u))\)
  \item \((\exists u)\,S(u) \rightarrow R(u,c)\)
  \item \((\exists u)\,(S(u) \rightarrow R(u,c))\)
\end{enumerate}
\ruangjawaban[8cm]

\soaluas
\textit{(Interpretasi \& Arti Kalimat)}\\
Diberikan ekspresi
\[
G :\;
(\forall u)\,(S(u) \rightarrow R(u,c)).
\]
\begin{enumerate}[label=\alph*., leftmargin=2em, itemsep=0.25em]
  \item Misalkan interpretasi \(I\) dengan domain \(D=\{1,2,3\}\),
        \(c_I = 2\),
        \(S_I = \{1,3\}\) (ekstensi predikat \(S\)),
        dan \(R_I = \{(1,2),(3,2),(2,2)\}\).
        Tuliskan arti \(G\) berdasarkan \(I\), lalu tentukan nilai kebenarannya.
        Tunjukkan pemeriksaan untuk setiap elemen domain.
  \item Misalkan interpretasi \(J\) dengan domain semua mata kuliah di kampus,
        \(c_J =\) ``Logika Matematika'',
        \(S_J(u)\): ``\(u\) adalah mata kuliah wajib'',
        \(R_J(u,v)\): ``\(u\) menjadi prasyarat bagi \(v\)''.
        Tuliskan arti \(G\) berdasarkan \(J\) dalam bahasa Indonesia.
\end{enumerate}
\ruangjawaban[9cm]

\soaluas
\textit{(Semantik Kuantor \& Urutan Kuantor)}\\
\begin{enumerate}[label=\alph*., leftmargin=2em, itemsep=0.25em]
  \item Domain \(D=\{\mathit{Andi}, \mathit{Dina}, \mathit{Citra}\}\).
        Predikat \(Loves(x,y)\) diinterpretasikan sebagai relasi
        \(\{(\mathit{Andi},\mathit{Dina}), (\mathit{Dina},\mathit{Citra}), (\mathit{Citra},\mathit{Citra})\}\).
        Tentukan nilai kebenaran:
        \begin{enumerate}[label=\roman*., itemsep=0.15em]
          \item \(\exists x\,\forall y\,Loves(x,y)\)
          \item \(\forall x\,\exists y\,Loves(x,y)\)
        \end{enumerate}
        Jelaskan singkat dengan menelusuri elemen domain.
  \item Misalkan \(Less(x,y)\) berarti ``\(x < y\)'' pada domain bilangan bulat.
        Bandingkan dua formula berikut:
        \[
        \forall x\,\exists y\,Less(x,y)
        \qquad\text{dan}\qquad
        \exists y\,\forall x\,Less(x,y).
        \]
        Manakah yang benar, manakah yang salah?
        Mengapa urutan kuantor penting?
\end{enumerate}
\ruangjawaban[9cm]

\soaluas
\textit{(Kecocokan / Simbol Bebas)}\\
Misalkan interpretasi \(I\) meliputi bilangan bulat dengan:
\(c \mapsto 0\), \(d \mapsto 5\), \(u \mapsto 2\),
\(h_I(n)=n+1\).
Interpretasi \(J\) meliputi bilangan bulat dengan:
\(c \mapsto 0\), \(u \mapsto 4\),
\(h_J(n)=n-1\).
\begin{enumerate}[label=\alph*., leftmargin=2em, itemsep=0.25em]
  \item Sebutkan simbol bebas dari ekspresi \(R(h(u),c)\).
  \item Apakah \(I\) dan \(J\) \textbf{cocok} untuk konstanta \(c\)? Untuk variabel \(u\)? Jelaskan.
  \item Apakah \(I\) dan \(J\) cocok untuk ekspresi \(h(c)\)? Untuk ekspresi \(h(d)\)? Jelaskan.
\end{enumerate}
\ruangjawaban[8cm]

\soaluas
\textit{(Validitas)}\\
Untuk setiap kalimat tertutup berikut, tentukan apakah \textbf{valid}.
Jika valid, berikan sketsa pembuktian singkat berdasarkan aturan semantik.
Jika tidak valid, berikan \textbf{satu} interpretasi yang membuat kalimat bernilai salah.
\begin{enumerate}[label=\alph*., leftmargin=2em, itemsep=0.25em]
  \item \(\forall u\,(S(u) \rightarrow S(u))\)
  \item \(\bigl(\exists u\,S(u)\bigr) \rightarrow \bigl(\forall u\,S(u)\bigr)\)
  \item \(\forall u\,\exists v\,R(u,v) \rightarrow \exists v\,\forall u\,R(u,v)\)
\end{enumerate}
\ruangjawaban[9cm]

\vspace{1em}
\begin{center}
  \textit{--- Selamat mengerjakan ---}
\end{center}

\end{document}
